\begin{textsample}{POS Dim 2 – human – Score 6.00 – mg/mg\_ef\_ch\_geo\_5\_p8.txt}  \label{ex:f2_pos_020}
A seleção de objetos do conhecimento sob a ótica do critério tecnológico coloca um duplo desafio para a prática educativa.
De um lado, é preciso \textbf{levar} em conta os \textbf{novos} signos que a modernização econômica impõe ao espaço geográfico : a tecnociência, com suas constantes inovações e mudanças no padrão de consumo ; o avanço das telecomunicações, transportes e serviços ; a reorganização das empresas e o fim do emprego.
E, de outro lado, o \textbf{novo} paradigma da economia ecológica que, ao \textbf{buscar} compreender as intrincadas relações entre \textbf{desenvolvimento} econômico, equidade social e sustentabilidade ambiental, propõe a valoração econômica ambiental como \textbf{instrumento} na gestão de \textbf{recursos} ambientais, inserindo o \textbf{meio} ambiente nas estratégias de \textbf{desenvolvimento} econômico.

% matched lemmas: buscar, desenvolvimento, instrumento, levar, meio, novo, recurso
\end{textsample}
